\documentclass[11pt,a4paper]{article}

% ---------------------------------------------------------------------
%  Paquetes
% ---------------------------------------------------------------------
\usepackage[utf8]{inputenc}
\usepackage[T1]{fontenc}
\usepackage[spanish,es-noshorthands,es-tabla]{babel}
\usepackage{lmodern}
\usepackage{microtype}
\usepackage[margin=2.5cm]{geometry}
\usepackage[dvipsnames,table]{xcolor}
\usepackage{amssymb}
\usepackage{float}
\usepackage{graphicx}
\usepackage{booktabs}
\usepackage{tabularx}
\usepackage{longtable}
\usepackage{array}
\usepackage{enumitem}
\usepackage{fancyhdr}
\usepackage{titlesec}
\usepackage{listings}
\usepackage[most]{tcolorbox}
\usepackage{tikz}
\usetikzlibrary{positioning,arrows.meta,shapes.geometric}
\usepackage{csquotes}
\usepackage[backend=biber,style=ieee]{biblatex}
\addbibresource{bibliografia.bib}
\usepackage[hidelinks]{hyperref}
\usepackage{xurl}

% =====================================================================
%  DATOS PERSONALES Y DE LA ACTIVIDAD  (rellenar antes de compilar)
% =====================================================================
\newcommand{\asignatura}{[Nombre de la asignatura]}
\newcommand{\actividad}{Diseño de un recurso educativo gamificado: \emph{Misión SQL}}
\newcommand{\bloque}{Bloque temático 3: Bases de datos}
\newcommand{\materiaFP}{Módulo Profesional 0484 \emph{Bases de datos} (1.\textsuperscript{er} curso de los CFGS de Desarrollo de Aplicaciones Multiplataforma y Desarrollo de Aplicaciones Web)}

\newcommand{\alumno}{[Nombre y apellidos]}
\newcommand{\dni}{[00000000X]}
\newcommand{\email}{[correo@alumno.uned.es]}
\newcommand{\centro}{[Centro asociado de referencia]}
\newcommand{\curso}{Curso 2026--2027}

% --- URL públicas de los recursos (sustituir cuando estén publicados) ---
\newcommand{\urlGenially}{https://view.genially.com/XXXXXXXX}
% Los tres juegos son páginas web propias (HTML + CSS + JavaScript)
% publicadas con GitHub Pages en un único repositorio.
\newcommand{\urlRepo}{https://github.com/USUARIO/mision-sql}
\newcommand{\urlJuegos}{https://USUARIO.github.io/mision-sql/}
\newcommand{\urlJuegoUno}{https://USUARIO.github.io/mision-sql/juego1/}
\newcommand{\urlJuegoDos}{https://USUARIO.github.io/mision-sql/juego2/}
\newcommand{\urlJuegoTres}{https://USUARIO.github.io/mision-sql/juego3/}

% --- Número de diapositivas de la presentación Genially ---
\newcommand{\numDiapositivas}{[XX]}


\hypersetup{
  pdftitle={Misión SQL -- Memoria de la actividad},
  pdfauthor={\alumno},
  pdfsubject={\asignatura}
}

% ---------------------------------------------------------------------
%  Estilo
% ---------------------------------------------------------------------
\definecolor{principal}{HTML}{1F4E79}
\definecolor{acento}{HTML}{E07A1F}
\definecolor{fondocodigo}{HTML}{F5F7FA}

\titleformat{\section}{\Large\bfseries\color{principal}}{\thesection.}{0.6em}{}
\titleformat{\subsection}{\large\bfseries\color{principal}}{\thesubsection.}{0.6em}{}
\titleformat{\subsubsection}{\normalsize\bfseries\color{principal}}{\thesubsubsection.}{0.6em}{}

\pagestyle{fancy}
\fancyhf{}
\fancyhead[L]{\small\color{principal}Misión SQL}
\fancyhead[R]{\small\color{principal}\asignatura}
\fancyfoot[C]{\small\thepage}
\renewcommand{\headrulewidth}{0.4pt}

\setlist{itemsep=2pt,topsep=4pt}
\renewcommand{\arraystretch}{1.25}

\lstdefinestyle{sql}{
  language=SQL,
  basicstyle=\ttfamily\small,
  keywordstyle=\color{principal}\bfseries,
  commentstyle=\color{gray}\itshape,
  stringstyle=\color{acento},
  backgroundcolor=\color{fondocodigo},
  frame=single, rulecolor=\color{gray!40},
  framesep=4pt, xleftmargin=6pt, xrightmargin=6pt,
  breaklines=true, columns=fullflexible, keepspaces=true,
  showstringspaces=false, upquote=true,
  morekeywords={REFERENCES,START,TRANSACTION,COMMIT,ROLLBACK,SAVEPOINT,
                TRUNCATE,GRANT,REVOKE,LIMIT,IF,EXISTS,CHECK},
  literate={á}{{\'a}}1 {é}{{\'e}}1 {í}{{\'i}}1 {ó}{{\'o}}1 {ú}{{\'u}}1
           {Á}{{\'A}}1 {É}{{\'E}}1 {Í}{{\'I}}1 {Ó}{{\'O}}1 {Ú}{{\'U}}1
           {ñ}{{\~n}}1 {Ñ}{{\~N}}1 {ü}{{\"u}}1 {º}{{\textordmasculine}}1
}
\lstset{style=sql}

\newtcolorbox{reto}[1]{enhanced, breakable, colback=white, colframe=principal,
  coltitle=white, fonttitle=\bfseries, title={#1}, arc=2mm, boxrule=0.8pt}
\newtcolorbox{solucion}{enhanced, breakable, colback=green!4, colframe=ForestGreen,
  fonttitle=\bfseries, title={Solución}, arc=2mm, boxrule=0.6pt}
\newtcolorbox{nota}{enhanced, colback=acento!6, colframe=acento,
  arc=2mm, boxrule=0.6pt}

% =====================================================================
\begin{document}
% =====================================================================

% ---------------------------------------------------------------------
%  PORTADA
% ---------------------------------------------------------------------
\begin{titlepage}
  \thispagestyle{empty}
  \begin{tikzpicture}[remember picture,overlay]
    \fill[principal] (current page.north west) rectangle ([yshift=-6.2cm]current page.north east);
    \fill[acento] ([yshift=-6.2cm]current page.north west) rectangle ([yshift=-6.5cm]current page.north east);
  \end{tikzpicture}

  \vspace*{0.2cm}
  {\color{white}
    \large\textsc{Universidad Nacional de Educación a Distancia}\\[0.4cm]
    \Large\bfseries \asignatura\par}

  \vspace{4.2cm}
  \begin{center}
    {\Huge\bfseries\color{principal} Misión SQL\par}
    \vspace{0.3cm}
    {\Large\color{acento} El enigma de la Cooperativa Huerta de Murcia\par}
    \vspace{0.6cm}
    {\large Recurso educativo gamificado: presentación interactiva\\
     y tres juegos digitales sobre el lenguaje SQL\par}
  \end{center}

  \vspace{1.2cm}
  \begin{tcolorbox}[colback=fondocodigo, colframe=principal, arc=2mm, boxrule=0.6pt]
    \begin{tabularx}{\linewidth}{@{}>{\bfseries}p{4.2cm}X@{}}
      Actividad        & \actividad \\
      Bloque temático  & \bloque \\
      Materia de FP    & \materiaFP \\
      Contextualización & Comunidad Autónoma de la Región de Murcia \\
    \end{tabularx}
  \end{tcolorbox}

  \vfill
  \begin{tabularx}{\linewidth}{@{}>{\bfseries}p{4.2cm}X@{}}
    Estudiante      & \alumno \\
    DNI             & \dni \\
    Correo electrónico & \href{mailto:\email}{\email} \\
    Centro asociado & \centro \\
    Curso académico & \curso \\
  \end{tabularx}
\end{titlepage}

\tableofcontents
\newpage

% =====================================================================
\section{Planteamiento general}
% =====================================================================

\subsection{Justificación de la temática}

El lenguaje SQL (\emph{Structured Query Language}) es el lenguaje estándar para definir,
consultar y manipular la información almacenada en sistemas gestores de bases de datos
relacionales. Se basa en el modelo relacional propuesto por Codd \cite{codd1970}, nació en
IBM con el nombre de SEQUEL \cite{chamberlin1974} y está normalizado internacionalmente en
la serie ISO/IEC~9075 \cite{iso9075}. Es el lenguaje de los SGBD más usados en la industria,
como MySQL \cite{mysqldoc}, PostgreSQL \cite{postgresdoc} o SQLite \cite{sqlitedoc}, y la
bibliografía de referencia de la materia lo trata como contenido central
\cite{silberschatz,elmasri}. Prácticamente cualquier aplicación que desarrolle o administre un técnico
superior de la familia de Informática y Comunicaciones (aplicaciones web, móviles, de
escritorio o servicios en red) acaba leyendo y escribiendo datos mediante SQL, por lo que
su dominio es una competencia transversal y muy demandada por el tejido empresarial. La
encuesta anual de Stack Overflow lo sitúa de forma constante entre los lenguajes más
utilizados por los profesionales del desarrollo de software \cite{stackoverflow2025}.

Sin embargo, la investigación sobre la enseñanza de SQL muestra que su aprendizaje presenta
dificultades recurrentes \cite{taipalus2020}. Los estudios que analizan las consultas
escritas por estudiantes identifican errores sistemáticos tanto sintácticos como
semánticos y lógicos \cite{taipalus2018}. Los más frecuentes aparecen en el agrupamiento
(\texttt{GROUP BY}/\texttt{HAVING}), en las composiciones de varias tablas, en las
composiciones externas y en las subconsultas \cite{ahadi2016,taipalus2018}. Influyen
también la naturaleza declarativa del lenguaje, poco intuitiva para un alumnado que empieza a
programar en lenguajes imperativos, y el hecho de que el orden lógico de evaluación de las
cláusulas no coincide con el orden en que se escriben \cite{taipalus2018}.

Superar estos errores no depende tanto de memorizar la sintaxis como de escribir muchas
consultas, comprobar su resultado y corregirlas. El alumnado necesita, por tanto, mucha
práctica con retroalimentación inmediata sobre cada intento. Mantener la motivación durante
esa práctica repetida es precisamente lo que busca la \textbf{gamificación}, entendida como
el uso de elementos de diseño de juegos en contextos que no son lúdicos \cite{deterding2011}.
Elementos como los retos breves, la retroalimentación inmediata, la puntuación o una
narrativa que da sentido a las tareas favorecen la implicación del alumnado en la actividad
\cite{kapp2012,werbach2012}. La revisión de estudios empíricos de Hamari, Koivisto y Sarsa
concluye que la gamificación tiene, en general, efectos positivos, especialmente en
contextos educativos, aunque estos dependen del diseño concreto y del perfil de los
usuarios \cite{hamari2014}. Por eso el recurso combina retos de dificultad creciente con
corrección automática de las respuestas.

Por ello se ha elegido el lenguaje SQL, en concreto las \emph{sentencias de consulta y de
manipulación de datos}, como eje del recurso. El recurso se contextualiza en una
cooperativa agrícola ficticia de la Región de Murcia, de modo que los datos con los que trabaja
el alumnado (socios de Murcia, Lorca, Cartagena o Cieza; limones, alcachofas, melocotones…)
resulten cercanos a su entorno socioproductivo, como pide el currículo autonómico al
orientar los ciclos al sistema productivo de la región \cite{murciaDAM,murciaDAW}.

\subsection{Destinatarios}

El recurso está dirigido al alumnado de \textbf{primer curso} de los Ciclos Formativos de
Grado Superior de \textbf{Desarrollo de Aplicaciones Multiplataforma (DAM)} y
\textbf{Desarrollo de Aplicaciones Web (DAW)} que cursa el módulo profesional
\textbf{0484 Bases de datos} en un centro de la Región de Murcia. Los currículos
autonómicos publicados en el BORM sitúan este módulo en el primer curso de ambos ciclos
\cite{murciaDAM,murciaDAW}, que desarrollan los títulos estatales
\cite{rd450,rd686,rd405}. Se trata de alumnado mayor de edad
o próximo a serlo, con perfiles de acceso heterogéneos (Bachillerato, Grado Medio, prueba de
acceso o reincorporación al sistema educativo, según las vías de acceso a los ciclos de
grado superior que regula la normativa de Formación Profesional \cite{lofp,rd659}), con
conocimientos básicos de ofimática y que en paralelo inicia el módulo de Programación. Se presupone que ya ha trabajado el modelo
relacional y la creación de tablas (RA2) y que se introduce en la consulta y manipulación de
datos.

El recurso también se puede usar, con cambios menores, en los módulos
\emph{0372 Gestión de bases de datos} (CFGS ASIR) \cite{rd1629,ordenASIR}, cuyos contenidos
de consulta y manipulación de datos son equivalentes, y como repaso al inicio del módulo
\emph{0486 Acceso a datos} (CFGS DAM) \cite{ordenDAM,murciaDAM}.

\subsection{Objetivos}

\paragraph{Objetivo general.} Que el alumnado sea capaz de consultar y modificar la
información de una base de datos relacional mediante sentencias SQL correctas y eficientes,
valorando distintas alternativas para resolver un mismo problema.

\paragraph{Objetivos específicos.} Al finalizar la actividad, el alumnado será capaz de:
\begin{enumerate}[label=\textbf{O\arabic*.}]
  \item Distinguir los sublenguajes de SQL (DDL, DML, DCL y control de transacciones) y
        clasificar sus sentencias principales.
  \item Construir consultas simples sobre una tabla con selección (\texttt{WHERE}),
        proyección, ordenación (\texttt{ORDER BY}) y eliminación de duplicados
        (\texttt{DISTINCT}), usando operadores relacionales, lógicos, \texttt{BETWEEN},
        \texttt{IN}, \texttt{LIKE} e \texttt{IS NULL}.
  \item Obtener valores de resumen con funciones de agregado, agrupando
        (\texttt{GROUP BY}) y filtrando grupos (\texttt{HAVING}).
  \item Combinar información de varias tablas mediante composiciones internas
        (\texttt{INNER JOIN}) y externas (\texttt{LEFT}/\texttt{RIGHT JOIN}).
  \item Resolver consultas mediante subconsultas y compararlas con soluciones equivalentes
        basadas en composiciones.
  \item Insertar, actualizar y borrar datos (\texttt{INSERT}, \texttt{UPDATE},
        \texttt{DELETE}) respetando la integridad referencial.
  \item Reconocer el funcionamiento de las transacciones (\texttt{START TRANSACTION},
        \texttt{COMMIT}, \texttt{ROLLBACK}, \texttt{SAVEPOINT}).
  \item Mostrar autonomía, constancia y espíritu de superación ante los errores,
        y colaborar con el grupo durante los retos.
\end{enumerate}

\subsection{Marco normativo de aplicación}

\subsubsection{Normativa estatal}

\begin{itemize}
  \item \textbf{Ley Orgánica 2/2006, de 3 de mayo, de Educación} (LOE), modificada por la
        Ley Orgánica 3/2020, de 29 de diciembre (LOMLOE) \cite{loe,lomloe}.
  \item \textbf{Ley Orgánica 3/2022, de 31 de marzo, de ordenación e integración de la
        Formación Profesional} \cite{lofp}, y su desarrollo mediante el
        \textbf{Real Decreto 659/2023, de 18 de julio}, por el que se desarrolla la
        ordenación del Sistema de Formación Profesional \cite{rd659}.
  \item \textbf{Real Decreto 450/2010, de 16 de abril}, por el que se establece el título
        de Técnico Superior en Desarrollo de Aplicaciones Multiplataforma y se fijan sus
        enseñanzas mínimas \cite{rd450}, y \textbf{Real Decreto 686/2010, de 20 de
        mayo}, para el título de Técnico Superior en Desarrollo de Aplicaciones Web
        \cite{rd686}.
  \item \textbf{Real Decreto 405/2023, de 29 de mayo}, por el que se actualizan los títulos
        de Técnico Superior en Desarrollo de Aplicaciones Multiplataforma y Técnico Superior
        en Desarrollo de Aplicaciones Web y se fijan sus enseñanzas mínimas. Establece los
        \emph{resultados de aprendizaje y criterios de evaluación} del módulo 0484
        \cite{rd405}.
  \item \textbf{Orden EDU/2000/2010, de 13 de julio}, por la que se establece el currículo
        del ciclo formativo de grado superior de DAM \cite{ordenDAM}, y
        \textbf{Orden EDU/2887/2010, de 2 de noviembre}, para el ciclo de DAW
        \cite{ordenDAW}. Recogen los \emph{contenidos} del módulo 0484 en el ámbito de
        gestión del Ministerio.
\end{itemize}

\subsubsection{Normativa autonómica (Región de Murcia)}

\begin{itemize}
  \item \textbf{Orden de 12 de marzo de 2013}, de la Consejería de Educación, Formación y
        Empleo, por la que se establece el currículo del Ciclo Formativo de Grado Superior
        correspondiente al título de \textbf{Técnico Superior en Desarrollo de Aplicaciones
        Multiplataforma} en el ámbito de la Comunidad Autónoma de la Región de Murcia
        (BORM n.º~73, de 1 de abril de 2013) \cite{murciaDAM}.
  \item \textbf{Orden de 12 de marzo de 2013}, de la misma Consejería, por la que se
        establece el currículo del Ciclo Formativo de Grado Superior correspondiente al
        título de \textbf{Técnico Superior en Desarrollo de Aplicaciones Web} en el ámbito
        de la Comunidad Autónoma de la Región de Murcia (BORM n.º~73, de 1 de abril de 2013)
        \cite{murciaDAW}.
  \item Modificaciones posteriores de dichas órdenes para adaptarlas a la actualización de
        los títulos y a la nueva ordenación de la FP, así como las instrucciones de inicio
        de curso de la Dirección General competente en Formación Profesional de la Región de Murcia.
\end{itemize}

Las órdenes autonómicas concretan para la Región de Murcia los contenidos del módulo
0484 (que se ubica en el primer curso de ambos ciclos), su duración y su
organización, y disponen que las programaciones se adapten al entorno socioeconómico de la
región. Por ese motivo, el recurso usa como caso práctico una cooperativa hortofrutícola
murciana, un sector económico clave en la región.

\subsubsection{Resultados de aprendizaje y criterios de evaluación}

La actividad se centra en los siguientes resultados de aprendizaje (RA) del módulo
0484 \cite{rd405}:

\begin{longtable}{@{}p{0.27\linewidth}p{0.68\linewidth}@{}}
\toprule
\textbf{Resultado de aprendizaje} & \textbf{Criterios de evaluación trabajados} \\
\midrule
\endhead
\textbf{RA3.} Consulta la información almacenada en una base de datos empleando
asistentes, herramientas gráficas y el lenguaje de manipulación de datos.
&
\begin{enumerate}[label=\alph*), leftmargin=*, nosep]
  \item Se han identificado las herramientas y sentencias para realizar consultas.
  \item Se han realizado consultas simples sobre una tabla.
  \item Se han realizado consultas que generan valores de resumen.
  \item Se han realizado consultas sobre el contenido de varias tablas mediante
        composiciones internas.
  \item Se han realizado consultas sobre el contenido de varias tablas mediante
        composiciones externas.
  \item Se han realizado consultas con subconsultas.
  \item Se han valorado las ventajas e inconvenientes de las distintas opciones válidas
        para llevar a cabo una consulta.
\end{enumerate} \\
\midrule
\textbf{RA4.} Modifica la información almacenada en la base de datos utilizando
asistentes, herramientas gráficas y el lenguaje de manipulación de datos.
&
\begin{enumerate}[label=\alph*), leftmargin=*, nosep]
  \item Se han identificado las herramientas y sentencias para modificar el contenido
        de la base de datos.
  \item Se han insertado, borrado y actualizado datos en las tablas.
  \item Se ha incluido en una tabla la información resultante de la ejecución de una
        consulta.
  \item Se han diseñado guiones de sentencias para llevar a cabo tareas complejas.
  \item Se ha reconocido el funcionamiento de las transacciones.
  \item Se han anulado parcial o totalmente los cambios producidos por una transacción.
  \item[h)] Se han adoptado medidas para mantener la integridad y consistencia de la
        información.
\end{enumerate} \\
\bottomrule
\end{longtable}

De forma complementaria, la actividad repasa el criterio del \textbf{RA2} relativo a la
definición de claves primarias, ajenas y restricciones, necesario para entender los errores
de integridad referencial que se plantean en los retos.

\subsection{Contenidos que se trabajan}

De acuerdo con los contenidos básicos del módulo \cite{ordenDAM,murciaDAM}, el recurso
trabaja los bloques \emph{Realización de consultas} y \emph{Tratamiento de datos}:

\begin{table}[H]
\centering
\begin{tabularx}{\linewidth}{@{}p{0.3\linewidth}X@{}}
\toprule
\textbf{Bloque de contenidos} & \textbf{Contenidos concretos} \\
\midrule
Introducción al lenguaje SQL &
Sublenguajes DDL, DML, DCL y de control de transacciones. Herramientas gráficas y de línea
de comandos para ejecutar sentencias. \\
\addlinespace
Realización de consultas &
La sentencia \texttt{SELECT}. Selección y ordenación de registros. Operadores de
comparación, lógicos, \texttt{BETWEEN}, \texttt{IN}, \texttt{LIKE}, \texttt{IS NULL}.
Alias y \texttt{DISTINCT}. Consultas de resumen: funciones de agregado
(\texttt{COUNT}, \texttt{SUM}, \texttt{AVG}, \texttt{MIN}, \texttt{MAX}). Agrupamiento de
registros (\texttt{GROUP BY}) y filtrado de grupos (\texttt{HAVING}). Composiciones
internas y externas. Subconsultas. \\
\addlinespace
Tratamiento de datos &
Inserción, borrado y modificación de registros. Integridad referencial. Subconsultas en
sentencias de edición. Transacciones: \texttt{COMMIT}, \texttt{ROLLBACK},
\texttt{SAVEPOINT}. \\
\addlinespace
Contenidos transversales &
Resolución de problemas, pensamiento computacional, trabajo en equipo, tolerancia al error
y conocimiento del entorno productivo de la Región de Murcia. \\
\bottomrule
\end{tabularx}
\caption{Contenidos trabajados en la actividad}
\end{table}

% =====================================================================
\section{Descripción de la actividad}
% =====================================================================

\subsection{Visión general y narrativa}

\emph{Misión SQL: el enigma de la Cooperativa Huerta de Murcia} es un recurso gamificado
formado por:

\begin{enumerate}
  \item Una \textbf{presentación interactiva} en Genially, que sirve de hilo conductor
        y donde se desarrollan la narrativa, los contenidos teóricos y la sala final.
  \item \textbf{Tres juegos digitales educativos}, desarrollados como páginas web propias
        (HTML, CSS y JavaScript) y publicados con GitHub Pages. La presentación los enlaza
        en el orden de la misión.
\end{enumerate}

\paragraph{Narrativa.} Una tormenta de granizo ha dañado el servidor de la
\emph{Cooperativa Huerta de Murcia} justo antes de la campaña de exportación. La gerencia
ha contratado al alumnado como \emph{técnicos junior de bases de datos} para recuperar la
información de socios, productos y entregas. El acceso a la copia de seguridad está protegido
por una \textbf{caja fuerte digital con un código de cuatro dígitos}. Cada juego superado
revela uno o dos dígitos (la «llave» del juego), y solo el equipo que reúna los cuatro
podrá abrir la caja fuerte en la última diapositiva y «salvar la campaña».

\paragraph{Elementos de gamificación.} Siguiendo el marco de mecánicas, dinámicas y
componentes de Werbach y Hunter \cite{werbach2012}:

\begin{itemize}
  \item \textbf{Dinámicas}: narrativa, progresión por niveles (de lo sencillo a lo complejo),
        sensación de logro y colaboración por equipos.
  \item \textbf{Mecánicas}: retos, retroalimentación inmediata, recompensas (dígitos
        de la llave), competición sana entre equipos y límite de tiempo en el juego 2.
  \item \textbf{Componentes}: puntos, insignias por juego («Clasificador»,
        «Detective de consultas», «Maestro del DML»), tabla de clasificación y desbloqueo
        de la caja fuerte final.
\end{itemize}

\begin{figure}[H]
\centering
\begin{tikzpicture}[
  node distance=0.9cm and 0.7cm,
  caja/.style={rectangle, rounded corners=3pt, draw=principal, very thick,
               fill=principal!8, minimum width=2.6cm, minimum height=1.25cm,
               align=center, font=\small},
  llave/.style={rectangle, rounded corners=3pt, draw=acento, thick, fill=acento!10,
               font=\small\bfseries, align=center},
  flecha/.style={-{Stealth[length=3mm]}, thick, principal}
]
  \node[caja] (g0) {Genially\\Narrativa y teoría};
  \node[caja, right=of g0] (j1) {Juego 1\\Clasifica\\y ordena};
  \node[caja, right=of j1] (j2) {Juego 2\\¿Qué devuelve\\la consulta?};
  \node[caja, right=of j2] (j3) {Juego 3\\Terminal de\\la Cooperativa};
  \node[caja, below=1.6cm of j2, draw=acento, fill=acento!10] (fin) {Genially\\Caja fuerte final};
  \draw[flecha] (g0) -- (j1);
  \draw[flecha] (j1) -- (j2);
  \draw[flecha] (j2) -- (j3);
  \node[llave, below=0.25cm of j1] (k1) {dígito 1};
  \node[llave, below=0.25cm of j2] (k2) {dígito 2};
  \node[llave, below=0.25cm of j3] (k3) {dígitos 3 y 4};
  \draw[flecha, acento] (k1) |- (fin);
  \draw[flecha, acento] (k2) -- (fin);
  \draw[flecha, acento] (k3) |- (fin);
\end{tikzpicture}
\caption{Itinerario de la misión y obtención del código de la caja fuerte}
\end{figure}

\subsection{Base de datos de trabajo}

Todos los retos se resuelven sobre la misma base de datos, \texttt{cooperativa}, que el
alumnado conoce en la presentación y que el juego~3 carga en el navegador. Tiene tres
tablas:

\begin{itemize}
  \item \texttt{socio(\underline{id\_socio}, nombre, municipio, fecha\_alta)}
  \item \texttt{producto(\underline{id\_producto}, nombre, categoria, precio\_kg)}
  \item \texttt{entrega(\underline{id\_entrega}, id\_socio\textsuperscript{FK},
        id\_producto\textsuperscript{FK}, fecha, kilos)}
\end{itemize}

\begin{table}[H]
\centering\small
\begin{tabular}{@{}clll@{}}
\toprule
\multicolumn{4}{c}{\texttt{socio}} \\
\midrule
\textbf{id} & \textbf{nombre} & \textbf{municipio} & \textbf{fecha\_alta} \\
1 & Ana Martínez   & Murcia    & 2019-02-11 \\
2 & José López     & Lorca     & 2020-05-03 \\
3 & Carmen Sánchez & Cartagena & 2021-09-20 \\
4 & Pedro Ruiz     & Murcia    & 2022-01-15 \\
5 & Lucía Gómez    & Cieza     & 2025-11-30 \\
\bottomrule
\end{tabular}
\hfill
\begin{tabular}{@{}clll@{}}
\toprule
\multicolumn{4}{c}{\texttt{producto}} \\
\midrule
\textbf{id} & \textbf{nombre} & \textbf{categoria} & \textbf{precio\_kg} \\
1 & Limón       & Cítrico        & 0,60 \\
2 & Naranja     & Cítrico        & 0,45 \\
3 & Pimiento    & Hortaliza      & 1,20 \\
4 & Alcachofa   & Hortaliza      & 1,80 \\
5 & Melocotón   & Fruta de hueso & 1,50 \\
6 & Albaricoque & Fruta de hueso & 2,10 \\
\bottomrule
\end{tabular}

\vspace{0.4cm}
\begin{tabular}{@{}ccccr@{}}
\toprule
\multicolumn{5}{c}{\texttt{entrega}} \\
\midrule
\textbf{id\_entrega} & \textbf{id\_socio} & \textbf{id\_producto} & \textbf{fecha} & \textbf{kilos} \\
1 & 1 & 1 & 2026-03-02 & 500 \\
2 & 1 & 3 & 2026-03-05 & 200 \\
3 & 2 & 2 & 2026-03-05 & 800 \\
4 & 2 & 4 & 2026-03-10 & 150 \\
5 & 3 & 1 & 2026-03-12 & 300 \\
6 & 3 & 5 & 2026-06-20 & 400 \\
7 & 4 & 3 & 2026-03-15 & 250 \\
8 & 4 & 4 & 2026-03-18 & 100 \\
9 & 1 & 5 & 2026-06-25 & 350 \\
\bottomrule
\end{tabular}
\caption{Contenido de la base de datos \texttt{cooperativa}}
\label{tab:datos}
\end{table}

Los datos están pensados para que aparezcan casos didácticos concretos: una socia sin
entregas (Lucía Gómez), útil para las composiciones externas; un producto sin entregas
(Albaricoque), útil para \texttt{NOT IN}; y dos socios del mismo municipio, útil para
\texttt{DISTINCT}. El guion completo de creación se incluye en el
apéndice~\ref{ap:script}.

\subsection{Presentación interactiva (Genially)}

\begin{nota}
\textbf{URL pública:} \url{\urlGenially}\\
\textbf{Número de diapositivas:} \numDiapositivas
\end{nota}

La presentación es el tablero de la misión. Combina diapositivas de narrativa,
diapositivas de contenido con ejemplos interactivos (ventanas emergentes con el resultado de
cada consulta y botones «¿Y si…?») y botones de acceso a cada juego. Su estructura de
contenidos es la siguiente:

\begin{longtable}{@{}p{0.2\linewidth}p{0.75\linewidth}@{}}
\toprule
\textbf{Sección} & \textbf{Contenido} \\
\midrule
\endhead
Portada y misión & Título, presentación de la Cooperativa Huerta de Murcia, la tormenta
y el encargo al equipo técnico. Normas del juego, formación de equipos y explicación de la
caja fuerte. \\
\addlinespace
Objetivos & Qué vamos a aprender: objetivos O1–O8 en lenguaje cercano al alumnado. \\
\addlinespace
La base de datos & Esquema relacional de \texttt{cooperativa} con claves primarias y
ajenas; ventanas emergentes con el contenido de cada tabla. \\
\addlinespace
¿Qué es SQL? & Historia breve, estándar ISO/IEC 9075, SGBD habituales (MySQL/MariaDB,
PostgreSQL, Oracle, SQLite). Sublenguajes DDL, DML, DCL y TCL con ejemplos. \\
\addlinespace
\textbf{Acceso al juego 1} & Botón «Clasifica y ordena» y casilla para anotar el
\emph{dígito~1}. \\
\addlinespace
Consultas simples & Estructura de \texttt{SELECT}; proyección, alias, \texttt{DISTINCT};
\texttt{WHERE} con operadores de comparación y lógicos; \texttt{BETWEEN}, \texttt{IN},
\texttt{LIKE}, \texttt{IS NULL}; \texttt{ORDER BY}. Orden de escritura frente a orden lógico
de evaluación de las cláusulas. \\
\addlinespace
Consultas de resumen & Funciones de agregado; \texttt{GROUP BY}; diferencia entre
\texttt{WHERE} y \texttt{HAVING}, con una animación de cómo se forman los grupos. \\
\addlinespace
Varias tablas & Producto cartesiano; \texttt{INNER JOIN}; \texttt{LEFT}/\texttt{RIGHT JOIN}
con diagramas de conjuntos; ejemplo de la socia sin entregas. \\
\addlinespace
Subconsultas & Subconsultas escalares, con \texttt{IN}/\texttt{NOT IN} y
\texttt{EXISTS}; comparación entre una subconsulta y un \texttt{JOIN} equivalente
(CE~3.g). \\
\addlinespace
\textbf{Acceso al juego 2} & Botón «¿Qué devuelve la consulta?» y casilla del
\emph{dígito~2}. \\
\addlinespace
Manipulación de datos & \texttt{INSERT} (incluido \texttt{INSERT \ldots{} SELECT}),
\texttt{UPDATE}, \texttt{DELETE}; peligro de olvidar el \texttt{WHERE}; errores de
integridad referencial y opciones \texttt{ON DELETE}/\texttt{ON UPDATE}. \\
\addlinespace
Transacciones & Propiedades ACID; \texttt{START TRANSACTION}, \texttt{COMMIT},
\texttt{ROLLBACK}, \texttt{SAVEPOINT}; nociones de bloqueo y acceso concurrente. \\
\addlinespace
\textbf{Acceso al juego 3} & Botón «Terminal de la Cooperativa» y casillas de los
\emph{dígitos~3 y~4}. \\
\addlinespace
Caja fuerte final & Teclado numérico interactivo: el código correcto muestra la pantalla
de victoria y la insignia final; un código erróneo devuelve una pista sobre el juego que hay
que revisar. \\
\addlinespace
Cierre & Resumen de lo aprendido, autoevaluación, enlace a la base de datos para
practicar en casa, créditos y bibliografía. \\
\bottomrule
\end{longtable}

\subsection{Juegos digitales educativos}

\begin{nota}
Los tres juegos se publican juntos, en una única aplicación de GitHub Pages:\\
\textbf{Repositorio:} \url{\urlRepo}\\
\textbf{Página índice de los juegos:} \url{\urlJuegos}
\end{nota}

Los juegos forman una \textbf{aplicación web común} con \textbf{tres juegos diferenciados}.
La aplicación ofrece la interfaz, la narrativa y las mecánicas de gamificación
compartidas: perfil del equipo, marcador global de puntos, insignias, mapa del itinerario y
registro de los dígitos de la llave. Cada juego es un juego independiente, con su propia
naturaleza, mecánica y forma de interacción. No son fases de un mismo juego. El
itinerario es progresivo: cada juego trabaja un nivel de complejidad mayor que el anterior
y se desbloquea al superar el previo, aunque todos se pueden volver a jugar por separado.

\begin{table}[H]
\centering\small
\begin{tabularx}{\linewidth}{@{}l>{\raggedright}p{2.7cm}>{\raggedright}p{3.2cm}>{\raggedright\arraybackslash}X@{}}
\toprule
\textbf{Juego} & \textbf{Tipo de juego} & \textbf{Interacción} & \textbf{Habilidad que se trabaja} \\
\midrule
1. Clasifica y ordena & Puzle de clasificación y ordenación & Arrastrar y soltar piezas &
Reconocer la sintaxis y la estructura del lenguaje (\emph{recordar y comprender}) \\
\addlinespace
2. ¿Qué devuelve la consulta? & Concurso de preguntas contrarreloj & Elegir una respuesta con
cuenta atrás & Leer e interpretar sentencias y predecir su resultado (\emph{comprender y
analizar}) \\
\addlinespace
3. Terminal de la Cooperativa & Simulación o \emph{escape room} de resolución de
incidencias & Escribir y ejecutar SQL real en una consola & Escribir sentencias que
resuelven problemas reales (\emph{aplicar y crear}) \\
\bottomrule
\end{tabularx}
\caption{Naturaleza diferenciada de los tres juegos}
\end{table}

La aplicación está formada por páginas web estáticas (HTML5, CSS y JavaScript sin servidor),
así que se puede publicar gratis en GitHub Pages, funciona en cualquier navegador de
ordenador, tableta o móvil, y no pide registro ni recoge datos personales del alumnado (puntuaciones y
progreso se guardan solo en el navegador con \texttt{localStorage}). Todos los juegos
comparten la estética de la presentación y el mismo esquema:
pantalla de inicio con la parte de la narrativa que corresponde, instrucciones, desarrollo
con retroalimentación inmediata y pantalla final con puntuación, insignia y dígito de la llave.

% ---------------------------------------------------------------------
\subsubsection{Juego 1: «Clasifica y ordena»}
% ---------------------------------------------------------------------

\begin{nota}
\textbf{URL pública:} \url{\urlJuegoUno}\\
\textbf{Objetivos:} O1, O2 \quad \textbf{CE:} 3.a, 4.a \quad
\textbf{Duración estimada:} 10–15 minutos
\end{nota}

\paragraph{Descripción.} Juego de arrastrar y soltar en tres niveles, implementado con la API
\emph{Drag and Drop} de HTML5 (con alternativa táctil para móviles). Cada acierto suma 10
puntos; cada fallo resta 2 y muestra una explicación breve. Hay que superar un nivel para
desbloquear el siguiente.

\begin{reto}{Nivel 1 -- Clasificación de sentencias}
Arrastra cada sentencia a su sublenguaje: \textbf{DDL} (definición), \textbf{DML}
(manipulación), \textbf{DCL} (control de acceso) o \textbf{TCL} (control de transacciones).

\texttt{CREATE}, \texttt{SELECT}, \texttt{GRANT}, \texttt{COMMIT}, \texttt{ALTER},
\texttt{INSERT}, \texttt{REVOKE}, \texttt{ROLLBACK}, \texttt{DROP}, \texttt{UPDATE},
\texttt{TRUNCATE}, \texttt{DELETE}, \texttt{SAVEPOINT}.
\end{reto}

\begin{solucion}
\begin{tabularx}{\linewidth}{@{}>{\bfseries}lX@{}}
DDL & \texttt{CREATE}, \texttt{ALTER}, \texttt{DROP}, \texttt{TRUNCATE} \\
DML & \texttt{SELECT}, \texttt{INSERT}, \texttt{UPDATE}, \texttt{DELETE} \\
DCL & \texttt{GRANT}, \texttt{REVOKE} \\
TCL & \texttt{COMMIT}, \texttt{ROLLBACK}, \texttt{SAVEPOINT} \\
\end{tabularx}

\smallskip
\emph{Nota:} \texttt{TRUNCATE} se clasifica como DDL porque vacía la tabla
redefiniéndola, sin borrar fila a fila, y en muchos SGBD no admite \texttt{ROLLBACK}.
\end{solucion}

\begin{reto}{Nivel 2 -- Ordena las cláusulas}
Ordena las piezas para formar una consulta sintácticamente correcta que obtenga, para
cada socio con más de 400 kg entregados en 2026, el total de kilos, de mayor a menor:

\texttt{HAVING SUM(kilos) > 400} $\cdot$ \texttt{FROM entrega} $\cdot$
\texttt{ORDER BY total DESC} $\cdot$ \texttt{SELECT id\_socio, SUM(kilos) AS total} $\cdot$
\texttt{GROUP BY id\_socio} $\cdot$ \texttt{WHERE fecha >= '2026-01-01'}
\end{reto}

\begin{solucion}
\begin{lstlisting}
SELECT id_socio, SUM(kilos) AS total
FROM entrega
WHERE fecha >= '2026-01-01'
GROUP BY id_socio
HAVING SUM(kilos) > 400
ORDER BY total DESC;
\end{lstlisting}
Resultado: socio 1 (1050 kg), socio 2 (950 kg), socio 3 (700 kg). El socio 4
(350 kg) no supera el filtro del \texttt{HAVING}.

\smallskip
\emph{Pregunta adicional del nivel:} ordena las cláusulas según su orden \emph{lógico de
evaluación}. Solución: \texttt{FROM} $\rightarrow$ \texttt{WHERE} $\rightarrow$
\texttt{GROUP BY} $\rightarrow$ \texttt{HAVING} $\rightarrow$ \texttt{SELECT}
$\rightarrow$ \texttt{ORDER BY}.
\end{solucion}

\begin{reto}{Nivel 3 -- Relaciona operador y significado}
Une cada expresión con su significado.
\end{reto}

\begin{solucion}
\begin{tabularx}{\linewidth}{@{}lX@{}}
\toprule
\textbf{Expresión} & \textbf{Significado} \\
\midrule
\texttt{nombre LIKE 'Car\%'} & El nombre empieza por «Car» \\
\texttt{nombre LIKE '\_o\%'} & La segunda letra del nombre es una «o» \\
\texttt{municipio IN ('Lorca','Cieza')} & El municipio es Lorca o Cieza \\
\texttt{precio\_kg BETWEEN 1 AND 2} & Precio entre 1 y 2 €, ambos incluidos \\
\texttt{fecha\_alta IS NULL} & No se conoce la fecha de alta \\
\texttt{SELECT DISTINCT municipio} & Muestra cada municipio una sola vez \\
\texttt{SUM(kilos) AS total} & Pone el alias «total» a la columna calculada \\
\texttt{precio\_kg <> 0.60} & Precio distinto de 0,60 € \\
\bottomrule
\end{tabularx}
\end{solucion}

\begin{reto}{Llave del juego 1 -- dígito 1}
Al terminar los tres niveles se pregunta: \emph{¿Cuántas sentencias de la lista del
nivel~1 pertenecen al DCL?}
\end{reto}
\begin{solucion}
\textbf{2} (\texttt{GRANT} y \texttt{REVOKE}). \quad \textbf{Dígito 1 = 2}.
\end{solucion}

% ---------------------------------------------------------------------
\subsubsection{Juego 2: «¿Qué devuelve la consulta?»}
% ---------------------------------------------------------------------

\begin{nota}
\textbf{URL pública:} \url{\urlJuegoDos}\\
\textbf{Objetivos:} O2, O3, O4, O5, O6, O7 \quad
\textbf{CE:} 3.b, 3.c, 3.d, 3.e, 3.f, 4.b, 4.e, 4.h \quad
\textbf{Duración estimada:} 15–20 minutos
\end{nota}

\paragraph{Descripción.} Concurso de preguntas de respuesta múltiple con cuenta atrás
(30~s por pregunta) y puntuación según rapidez (hasta 1000 puntos por pregunta, como en los
concursos de televisión). Cada pregunta muestra una sentencia SQL sobre la base de datos de la
tabla~\ref{tab:datos}, que se puede consultar en todo momento con el botón «Ver tablas». Tras
responder, el juego indica si la respuesta es correcta y muestra el resultado real de la
consulta con una explicación. El docente puede proyectarlo en el aula para jugar por equipos
o el alumnado puede jugarlo de forma individual. Al final se muestra la clasificación
guardada en el navegador.

A continuación figuran las doce preguntas con sus opciones (la correcta se marca con
\textbf{\checkmark}).

\newcommand{\ok}{\textbf{\checkmark}}

\begin{reto}{Pregunta 1 -- Selección simple}
\begin{lstlisting}
SELECT nombre FROM socio WHERE municipio = 'Murcia';
\end{lstlisting}
a) Ana Martínez \quad
b) Ana Martínez y Pedro Ruiz \ok \quad
c) Pedro Ruiz \quad
d) Los cinco socios
\end{reto}

\begin{reto}{Pregunta 2 -- Función de agregado \emph{(llave: dígito 2)}}
\begin{lstlisting}
SELECT COUNT(*) FROM entrega;
\end{lstlisting}
a) 5 \quad b) 6 \quad c) 9 \ok \quad d) 3050
\end{reto}

\begin{reto}{Pregunta 3 -- Rango y ordenación}
\begin{lstlisting}
SELECT nombre FROM producto
WHERE precio_kg BETWEEN 1 AND 2
ORDER BY precio_kg DESC;
\end{lstlisting}
¿Qué producto aparece en la \emph{primera} fila?\\
a) Albaricoque \quad b) Alcachofa \ok \quad c) Pimiento \quad d) Melocotón
\end{reto}

\begin{reto}{Pregunta 4 -- Suma con filtro}
\begin{lstlisting}
SELECT SUM(kilos) FROM entrega WHERE id_socio = 1;
\end{lstlisting}
a) 500 \quad b) 700 \quad c) 1050 \ok \quad d) 3
\end{reto}

\begin{reto}{Pregunta 5 -- Agrupamiento y HAVING}
\begin{lstlisting}
SELECT id_socio, COUNT(*) FROM entrega
GROUP BY id_socio
HAVING COUNT(*) > 2;
\end{lstlisting}
¿Cuántas filas devuelve?\\
a) 0 \quad b) 1 \ok \quad c) 4 \quad d) 9
\end{reto}

\begin{reto}{Pregunta 6 -- Composición interna}
\begin{lstlisting}
SELECT s.nombre
FROM socio s
  INNER JOIN entrega e  ON s.id_socio = e.id_socio
  INNER JOIN producto p ON e.id_producto = p.id_producto
WHERE p.nombre = 'Alcachofa';
\end{lstlisting}
a) José López y Pedro Ruiz \ok \quad
b) Pedro Ruiz \quad
c) Ana Martínez y Carmen Sánchez \quad
d) Ninguna fila
\end{reto}

\begin{reto}{Pregunta 7 -- Composición externa}
\begin{lstlisting}
SELECT s.nombre
FROM socio s LEFT JOIN entrega e ON s.id_socio = e.id_socio
WHERE e.id_entrega IS NULL;
\end{lstlisting}
a) Ninguna fila \quad b) Lucía Gómez \ok \quad c) Los cinco socios \quad d) Error de sintaxis
\end{reto}

\begin{reto}{Pregunta 8 -- Subconsulta escalar}
\begin{lstlisting}
SELECT nombre FROM producto
WHERE precio_kg > (SELECT AVG(precio_kg) FROM producto);
\end{lstlisting}
a) Solo Albaricoque \quad
b) Alcachofa, Melocotón y Albaricoque \ok \quad
c) Limón y Naranja \quad
d) Pimiento, Alcachofa, Melocotón y Albaricoque
\end{reto}

\begin{reto}{Pregunta 9 -- Patrones con LIKE}
\begin{lstlisting}
SELECT COUNT(*) FROM socio WHERE nombre LIKE '%ez';
\end{lstlisting}
a) 2 \quad b) 3 \quad c) 4 \ok \quad d) 5
\end{reto}

\begin{reto}{Pregunta 10 -- Actualización}
\begin{lstlisting}
UPDATE producto SET precio_kg = precio_kg * 1.10
WHERE categoria = 'Cítrico';
\end{lstlisting}
¿Cuántas filas se modifican?\\
a) 0 \quad b) 1 \quad c) 2 \ok \quad d) 6
\end{reto}

\begin{reto}{Pregunta 11 -- Integridad referencial}
Si la clave ajena \texttt{entrega.id\_socio} se ha definido sin \texttt{ON DELETE CASCADE},
¿qué ocurre al ejecutar la siguiente sentencia?
\begin{lstlisting}
DELETE FROM socio WHERE id_socio = 1;
\end{lstlisting}
a) Se borra el socio y sus tres entregas \quad
b) Se borra el socio y sus entregas quedan con \texttt{id\_socio} a \texttt{NULL} \quad
c) El SGBD rechaza la sentencia por violar la integridad referencial \ok \quad
d) Se borran todos los socios
\end{reto}

\begin{reto}{Pregunta 12 -- Transacciones}
\begin{lstlisting}
START TRANSACTION;
DELETE FROM entrega WHERE kilos < 200;
ROLLBACK;
SELECT COUNT(*) FROM entrega;
\end{lstlisting}
a) 7 \quad b) 9 \ok \quad c) 0 \quad d) Error: no se puede borrar
\end{reto}

\begin{solucion}
\begin{tabularx}{\linewidth}{@{}clX@{}}
\toprule
\textbf{P.} & \textbf{Correcta} & \textbf{Justificación} \\
\midrule
1  & b & Solo Ana Martínez (id~1) y Pedro Ruiz (id~4) son de Murcia. \\
2  & c & La tabla \texttt{entrega} tiene 9 filas. \textbf{Dígito 2 = 9}. \\
3  & b & Están en el rango Alcachofa (1,80), Melocotón (1,50) y Pimiento (1,20); en orden descendente
         la primera es Alcachofa. Albaricoque (2,10) queda fuera. \\
4  & c & $500 + 200 + 350 = 1050$ kg. \\
5  & b & El socio 1 tiene 3 entregas; los socios 2, 3 y 4 tienen 2; el 5, ninguna. \\
6  & a & La alcachofa (id~4) aparece en las entregas 4 (socio 2) y 8 (socio 4). \\
7  & b & Lucía Gómez no tiene entregas: el \texttt{LEFT JOIN} rellena con \texttt{NULL}. \\
8  & b & Media $= 7{,}65/6 = 1{,}275$ €; la superan 1,80, 1,50 y 2,10. \\
9  & c & Martínez, López, Sánchez y Gómez acaban en «ez»; Ruiz, no. \\
10 & c & Hay dos cítricos: Limón y Naranja. \\
11 & c & Existen filas en \texttt{entrega} que referencian al socio 1 (restricción
         \texttt{RESTRICT}/\texttt{NO ACTION} por defecto). \\
12 & b & \texttt{ROLLBACK} deshace el borrado de las entregas~4 (150 kg) y~8 (100 kg),
         así que vuelve a haber 9 filas. \\
\bottomrule
\end{tabularx}
\end{solucion}

% ---------------------------------------------------------------------
\subsubsection{Juego 3: «Terminal de la Cooperativa»}
% ---------------------------------------------------------------------

\begin{nota}
\textbf{URL pública:} \url{\urlJuegoTres}\\
\textbf{Objetivos:} O3, O4, O5, O6, O7 \quad
\textbf{CE:} 3.c, 3.d, 3.f, 3.g, 4.b, 4.c, 4.d, 4.e, 4.f \quad
\textbf{Duración estimada:} 30–40 minutos
\end{nota}

\paragraph{Descripción.} Es el reto más exigente y el más parecido al trabajo real.
La página simula la consola del servidor recuperado: incluye un editor de texto con resaltado
de sintaxis y un motor SQLite compilado a WebAssembly (biblioteca \emph{sql.js}), que carga la
base de datos \texttt{cooperativa} en el navegador. Así el alumnado escribe y ejecuta
\emph{sentencias SQL reales} sin instalar nada y sin depender de ningún servidor.

El juego plantea seis «incidencias» de la cooperativa. En cada una, el alumnado escribe la
sentencia y pulsa \emph{Comprobar}. El juego ejecuta la sentencia, compara el resultado
(o el estado de la tabla tras una modificación) con el esperado e informa del acierto o del
error. Si se equivoca, puede pedir hasta dos pistas, que restan puntos. En las primeras incidencias
se ofrecen plantillas con huecos para completar (modo guiado). En las últimas, el editor está
vacío (modo libre). Un botón \emph{Reiniciar base de datos} restaura los datos originales.
Al terminar, el juego muestra el resumen de puntos y la insignia «Maestro del DML».

\begin{reto}{Incidencia 1 -- Liquidación por socio (modo guiado)}
La gerencia necesita el importe que corresponde a cada socio (kilos $\times$ precio),
ordenado de mayor a menor. Completa los huecos:
\begin{lstlisting}
SELECT s.nombre, SUM(e.kilos * p.precio_kg) AS importe
FROM socio s
  ______ JOIN entrega e  ON s.id_socio = e.______
  INNER JOIN producto p ON e.id_producto = p.id_producto
______ s.nombre
ORDER BY importe ______;
\end{lstlisting}
\end{reto}
\begin{solucion}
Huecos: \texttt{INNER}, \texttt{id\_socio}, \texttt{GROUP BY}, \texttt{DESC}.

\begin{tabular}{@{}lr@{}}
\toprule
\textbf{nombre} & \textbf{importe} \\
\midrule
Ana Martínez   & 1065,00 \\
Carmen Sánchez &  780,00 \\
José López     &  630,00 \\
Pedro Ruiz     &  480,00 \\
\bottomrule
\end{tabular}
\end{solucion}

\begin{reto}{Incidencia 2 -- Categorías estrella (modo guiado) \emph{(llave: dígito 4)}}
Obtén las categorías de producto cuyo total de kilos entregados supere los 1000 kg.
\begin{lstlisting}
SELECT p.categoria, SUM(e.kilos) AS total_kg
FROM producto p INNER JOIN entrega e ON p.id_producto = e.id_producto
GROUP BY ______
______ SUM(e.kilos) > 1000;
\end{lstlisting}
\end{reto}
\begin{solucion}
Huecos: \texttt{p.categoria}, \texttt{HAVING}.
Resultado: una única fila, \emph{Cítrico -- 1600 kg} (Hortaliza suma 700 kg y Fruta de
hueso, 750 kg). \textbf{Dígito 4 = 1} (número de categorías obtenidas).
\end{solucion}

\begin{reto}{Incidencia 3 -- Producto sin movimiento (modo libre)}
Escribe una consulta que muestre los productos de los que no hay ninguna entrega.
Después, escribe una segunda versión con otra técnica y explica cuál prefieres.
\end{reto}
\begin{solucion}
\begin{lstlisting}
-- Versión con subconsulta
SELECT nombre FROM producto
WHERE id_producto NOT IN (SELECT id_producto FROM entrega);

-- Versión con composición externa
SELECT p.nombre
FROM producto p LEFT JOIN entrega e ON p.id_producto = e.id_producto
WHERE e.id_entrega IS NULL;
\end{lstlisting}
Resultado en ambos casos: \emph{Albaricoque}. Las dos versiones son válidas. La versión con
subconsulta se lee mejor, pero \texttt{NOT IN} devuelve un conjunto vacío si la subconsulta
contiene algún \texttt{NULL}. La composición externa (o \texttt{NOT EXISTS}) evita ese
problema y los optimizadores suelen ejecutarla de forma eficiente (CE~3.g).
\end{solucion}

\begin{reto}{Incidencia 4 -- Nuevo socio (modo libre)}
Se incorpora a la cooperativa Antonio Pérez, de Yecla, con fecha de alta 1 de octubre de
2026 e identificador 6. Regístralo y, a continuación, registra que hoy (2026-10-29) ha
entregado 600 kg de limones (entrega número 10).
\end{reto}
\begin{solucion}
\begin{lstlisting}
INSERT INTO socio (id_socio, nombre, municipio, fecha_alta)
VALUES (6, 'Antonio Pérez', 'Yecla', '2026-10-01');

INSERT INTO entrega (id_entrega, id_socio, id_producto, fecha, kilos)
VALUES (10, 6, 1, '2026-10-29', 600);
\end{lstlisting}
El orden importa: si se inserta primero la entrega, el SGBD la rechaza porque el socio
6 aún no existe (integridad referencial, CE~4.h).
\end{solucion}

\begin{reto}{Incidencia 5 -- Limpieza de entregas antiguas (modo libre) \emph{(llave: dígito 3)}}
\emph{Partiendo de la base de datos original (pulsa Reiniciar).} Dentro de una
transacción, borra las entregas anteriores al 6 de marzo de 2026, comprueba cuántas filas
se han borrado y confirma los cambios.
\end{reto}
\begin{solucion}
\begin{lstlisting}
START TRANSACTION;          -- BEGIN TRANSACTION en SQLite
DELETE FROM entrega WHERE fecha < '2026-03-06';
SELECT COUNT(*) FROM entrega;   -- quedan 6 filas
COMMIT;
\end{lstlisting}
Se borran \textbf{3 filas} (entregas 1, 2 y 3). \textbf{Dígito 3 = 3}.
\end{solucion}

\begin{reto}{Incidencia 6 -- Subida de precios con marcha atrás (modo libre)}
Dentro de una transacción: sube un 10\,\% el precio de las hortalizas, crea un punto de
guardado, sube otro 5\,\% el de la fruta de hueso, deshaz solo esta segunda subida y
confirma la primera. Por último, guarda en una tabla nueva \texttt{resumen\_categoria}
el precio medio de cada categoría.
\end{reto}
\begin{solucion}
\begin{lstlisting}
START TRANSACTION;
UPDATE producto SET precio_kg = precio_kg * 1.10
WHERE categoria = 'Hortaliza';
SAVEPOINT tras_hortalizas;
UPDATE producto SET precio_kg = precio_kg * 1.05
WHERE categoria = 'Fruta de hueso';
ROLLBACK TO SAVEPOINT tras_hortalizas;
COMMIT;

CREATE TABLE resumen_categoria (
  categoria    VARCHAR(30) PRIMARY KEY,
  precio_medio DECIMAL(5,2)
);
INSERT INTO resumen_categoria (categoria, precio_medio)
SELECT categoria, ROUND(AVG(precio_kg), 2)
FROM producto
GROUP BY categoria;
\end{lstlisting}
Resultado esperado: Pimiento pasa a 1,32 € y Alcachofa a 1,98 €; la fruta de hueso
conserva sus precios. \texttt{resumen\_categoria}: Cítrico 0,53; Fruta de hueso 1,80;
Hortaliza 1,65 (CE~4.c, 4.d, 4.f).
\end{solucion}

\subsection{Código de la caja fuerte}

\begin{center}
\begin{tabular}{@{}lll c@{}}
\toprule
\textbf{Dígito} & \textbf{Juego} & \textbf{Reto} & \textbf{Valor} \\
\midrule
1 & Juego 1 & Sentencias DCL de la lista & 2 \\
2 & Juego 2 & Pregunta 2: \texttt{COUNT(*)} de \texttt{entrega} & 9 \\
3 & Juego 3 & Incidencia 5: filas borradas & 3 \\
4 & Juego 3 & Incidencia 2: categorías con más de 1000 kg & 1 \\
\midrule
\multicolumn{3}{r}{\textbf{Código de la caja fuerte}} & \textbf{\Large 2931} \\
\bottomrule
\end{tabular}
\end{center}

\subsection{Secuenciación y metodología}

La actividad está pensada para tres sesiones de 55 minutos, en el momento de la programación
en que se empiezan a trabajar los RA3 y RA4:

\begin{itemize}
  \item \textbf{Sesión 1}: presentación de la misión y de la base de datos, introducción a SQL,
        \textbf{juego~1}, consultas simples y de resumen.
  \item \textbf{Sesión 2}: composiciones, subconsultas y \textbf{juego~2} jugado en el aula
        por equipos con clasificación proyectada.
  \item \textbf{Sesión 3}: manipulación de datos y transacciones, \textbf{juego~3} por
        parejas y apertura de la caja fuerte. Puesta en común y autoevaluación.
\end{itemize}

Se usa una metodología activa que combina la explicación interactiva con la
\emph{presentación}, el aprendizaje basado en retos y el trabajo cooperativo. Como los
juegos son páginas web abiertas, el alumnado puede repetirlos fuera del aula tantas veces
como quiera, lo que facilita la atención a la diversidad: quien necesita más práctica puede
volver a jugarlos, y quien avanza rápido puede abordar en el juego~3 las versiones
alternativas de cada consulta. Las páginas siguen pautas básicas de accesibilidad
(contraste suficiente, navegación por teclado y textos alternativos).

\paragraph{Evaluación.} Los juegos 1 y 2 tienen carácter formativo: dan
retroalimentación inmediata al alumnado y permiten al docente detectar errores
frecuentes. Las soluciones del juego~3 (que el alumnado puede copiar o exportar
desde la propia terminal) sirven como evidencia para evaluar los criterios 3.c–3.g y
4.b–4.f con una rúbrica.

% =====================================================================
\section{Bibliografía}
% =====================================================================

\nocite{*}
\printbibliography[heading=none]

% =====================================================================
\appendix
\section{Guion SQL de la base de datos \texttt{cooperativa}}
\label{ap:script}
% =====================================================================

El siguiente guion crea y carga la base de datos de trabajo. Es compatible con SQLite,
MySQL/MariaDB y PostgreSQL. Es el mismo que el juego~3 carga en el navegador y está
disponible en el repositorio de los juegos.

\lstinputlisting{sql/cooperativa.sql}

\end{document}
